\uuid{XO40}
\exo7id{3521}
\titre{exo7 3521}
\auteur{quercia}
\organisation{exo7}
\datecreate{2010-03-10}
\isIndication{false}
\isCorrection{true}
\chapitre{Réduction d'endomorphisme, polynôme annulateur}
\sousChapitre{Autre}
\cpge{mpsi-mp2i.limites-continuite.b, pcsi.limites-continuite.b, ptsi.limites-continuite.b, mpsi-mp2i.limites-continuite.c, pcsi.limites-continuite.c, ptsi.limites-continuite.c, mpsi-mp2i.applications-lineaires.b, pcsi.applications-lineaires.b, ptsi.applications-lineaires.b, mpsi-mp2i.integration.c, pcsi.integration.b, ptsi.integration.b, bcpst1.limites-continuite.c, mp-mpi.reduction-endomorphismes.b, pc.reduction.a, psi.reduction-endomorphismes.a, pt.reduction-endomorphismes.a, bcpst2.valeurs-propres-vecteurs-propres.a}
\module{Algèbre}
\niveau{L2}
\difficulte{}

\contenu{
\texte{
Soit $E = \mathcal{C}([0,+\infty[ \to\R)$
et $u : E \to  E, f \mapsto {\tilde f}$
avec $\tilde f(x) = \frac1x \int_{t=0}^x f(t)\,d t$.
}
\begin{enumerate}
    \item \question{Montrer que $\tilde f$ peut être prolongée en une fonction continue sur
    $[0,+\infty[$.}
    \item \question{Chercher les valeurs propres et les vecteurs propres de $u$.}
\reponse{
2. $0 < \lambda \le 1$ : $f(x) = Cx^{1/\lambda-1}$.
}
\end{enumerate}
}
